\documentclass[11pt,a4paper]{article}
\usepackage[utf8]{inputenc}
\usepackage[T1]{fontenc}
\usepackage{geometry}
\usepackage{amsmath}
\usepackage{amssymb}
\usepackage{array}
\usepackage{hyperref}

\geometry{margin=2cm, top=2.2cm, bottom=2.2cm}
\hypersetup{colorlinks=true, linkcolor=blue, urlcolor=blue}

\title{\textbf{Optimisation hivernale du d\'eneigement de Montr\'eal}\\[0.4em]
       \large Rapport ERO1 --- Formalisation et analyse}
\author{Rayan Kheroua \and Rayan Tail \and Arsan Abdi \and Yazid Tarmoul \and Lucas Defaud\\[0.3em]
        \normalsize Groupe 9 --- APPING1 --- EPITA}
\date{Juin 2026}

\begin{document}
\maketitle
\tableofcontents
\newpage

%=============================================================
\section{Formalisation du probl\`eme et m\'ethode de r\'esolution}
%=============================================================

\subsection{Donn\'ees et p\'erim\`etre}

Les donn\'ees utilis\'ees proviennent d'OpenStreetMap (OSM), t\'el\'echarg\'ees
via la biblioth\`eque Python OSMnx~\cite{osmnx}. Pour chaque arrondissement,
on extrait le r\'eseau routier carrossable (\texttt{network\_type="drive"}),
simplifi\'e topologiquement et projet\'e en coordonn\'ees UTM (distances en
m\`etres exactes). Les param\`etres \'economiques fournis sont :

\begin{itemize}
  \item Co\^ut fixe par v\'ehicule et par jour : $500$\,\$/jour
  \item Co\^ut kilom\'etrique : $1{,}1$\,\$/km
  \item Co\^ut horaire (8 premi\`eres heures) : $1{,}1$\,\$/h
  \item Co\^ut horaire (au-del\`a de 8\,h) : $1{,}3$\,\$/h
  \item Vitesse moyenne : $10$\,km/h
\end{itemize}

Les quatre arrondissements \'etudi\'es sont pr\'esent\'es dans le Tableau~1.

\begin{center}
\textbf{Tableau 1} --- R\'eseaux routiers des arrondissements\\[0.5em]
\begin{tabular}{|l|r|r|r|}
\hline
\textbf{Arrondissement} & \textbf{N\oe uds} & \textbf{Arcs} & \textbf{Longueur (km)} \\
\hline
Outremont                    &    236 &     598 &    71,1 \\
Verdun                       &    305 &     779 &   102,2 \\
Anjou                        &    737 &   1\,732 &   231,2 \\
Rivi\`ere-des-Prairies--PAT  &  1\,927 &   5\,215 &   715,8 \\
\hline
\end{tabular}
\end{center}

\medskip
\noindent\textbf{Contraintes prises en compte :}
sens de circulation (graphe dirig\'e) ;
obligation de parcourir chaque rue au moins une fois ;
circuit ferm\'e.

\noindent\textbf{Contraintes non prises en compte :}
capacit\'e des engins, horaires de collecte, partage multi-v\'ehicules (VRP).

\subsection{Hypoth\`eses et formalisation}

Le r\'eseau routier est model\'e comme un graphe dirig\'e pond\'er\'e
$G = (V,\, E,\, \ell)$ o\`u :
\begin{itemize}
  \item $V$ = ensemble des intersections (n\oe uds OSM),
  \item $E \subseteq V \times V$ = ensemble des tron\c{c}ons routiers (arcs),
  \item $\ell : E \to \mathbb{R}_+$ = longueur en m\`etres de chaque arc.
\end{itemize}

\noindent\textbf{Variable de d\'ecision :}
$x_e \in \mathbb{N},\; x_e \geq 1$ pour tout $e \in E$ (nombre de passages
sur l'arc $e$).

\noindent\textbf{Fonction objectif :}
\begin{equation}
  \min \sum_{e \in E} x_e \cdot \ell(e)
\end{equation}

\noindent\textbf{Contrainte d'eul\'erianit\'e :} pour que le circuit soit ferm\'e,
\begin{equation}
  \forall v \in V : \sum_{(v,w)\in E} x_{(v,w)} \;=\; \sum_{(u,v)\in E} x_{(u,v)}
\end{equation}

Ce probl\`eme est le \textbf{Probl\`eme du Postier Chinois Dirig\'e (PPC-D)},
r\'esolu en temps polynomial.

\subsection{M\'ethode de r\'esolution}

La r\'esolution se d\'eroule en trois \'etapes :

\begin{enumerate}
  \item \textbf{Extraction de la composante fortement connexe (CFC) principale.}
    Un circuit eul\'erien n'existe que sur un graphe fortement connexe.
    On retient la plus grande CFC ($\geq 97\,\%$ du r\'eseau dans tous les cas).

  \item \textbf{Calcul des d\'es\'equilibres et flot \`a co\^ut minimum.}
    Pour chaque n\oe ud $v$, on d\'efinit le d\'es\'equilibre
    $b_v = \deg^+(v) - \deg^-(v)$ (calculé sur le multigraphe original).
    Les n\oe uds avec $b_v > 0$ sont des \textit{puits} ;
    ceux avec $b_v < 0$ sont des \textit{sources}.
    On construit un r\'eseau de flot $F$ reliant chaque source \`a chaque puits,
    avec un co\^ut \'egal au plus court chemin dans le graphe simplifi\'e $H$.
    La r\'esolution via \texttt{nx.min\_cost\_flow} (algorithme des plus courts
    chemins successifs) donne le multiensemble d'arcs de \textit{deadhead}
    (trajets \`a vide) \`a ajouter pour \'equilibrer le graphe.

  \item \textbf{Circuit eul\'erien.}
    Sur le graphe augment\'e, le circuit est calcul\'e par l'algorithme de
    Hierholzer en $O(|E|)$.
\end{enumerate}

\noindent\textbf{Complexit\'e.}
La phase de flot est en $O(|V|^2 \log|V| \cdot |E|)$ dans le pire cas,
tractable pour les tailles consid\'er\'ees.

\medskip
\noindent\textbf{Indicateurs g\'en\'eriques :}
\begin{itemize}
  \item Distance totale du circuit $D_{\text{tot}}$ (km)
  \item Distance de deadhead $D_{\text{vide}}$ (km) et ratio $D_{\text{vide}}/D_{\text{tot}}$
  \item Taux de couverture $C = |\text{arcs parcourus}| / |E| \times 100\,\%$
  \item Co\^ut total $C_{\text{tot}}$ (\$) selon le mod\`ele de co\^ut
\end{itemize}

\subsection{Limites du mod\`ele}

\begin{itemize}
  \item \textbf{CFC partielle} : les arcs hors de la CFC (culs-de-sac) ne sont
    pas couverts, d'o\`u une couverture l\'eg\`erement inf\'erieure \`a 100\,\%.
  \item \textbf{Un seul v\'ehicule par circuit} : le partage entre v\'ehicules
    introduit un probl\`eme VRP NP-difficile.
  \item \textbf{Vitesse constante} : trafic, feux et densit\'e de neige ignor\'es.
  \item \textbf{Priorisation bas\'ee sur le type de route} uniquement (pas de
    donn\'ees POI sur les h\^opitaux ou commerces).
\end{itemize}

%=============================================================
\section{D\'efinition des trois sc\'enarios de priorisation}
%=============================================================

\subsection{Sc\'enario 1 --- Urgences et s\'ecurit\'e publique}

\noindent\textbf{Argumentaire.}
Lors d'une temp\^ete, l'acc\`es des v\'ehicules d'urgence est prioritaire.
Selon Sant\'e Montr\'eal~\cite{sante}, les d\'elais d'intervention augmentent de
40\,\% lors d'\'episodes neigeux. Les art\`eres primaires et secondaires
constituent le r\'eseau d'acc\`es aux h\^opitaux et casernes.

\noindent\textbf{D\'efinition des priorit\'es (attribut OSM \texttt{highway}) :}
\begin{itemize}
  \item Priorit\'e 3 (haute) : \texttt{primary}, \texttt{secondary}
  \item Priorit\'e 2 (moyenne) : \texttt{tertiary}
  \item Priorit\'e 1 (basse) : \texttt{residential}, \texttt{service}
\end{itemize}

\noindent\textbf{B\'en\'efices attendus :}
r\'eduction du d\'elai d'intervention des secours ;
maintien de la cha\^ine de soins.
\textbf{Cibles :} services d'urgence, patients.
\textbf{Risques :} les rues r\'esidentielles sont d\'egag\'ees en dernier.

\noindent\textbf{Indicateurs sp\'ecifiques :}
temps pour d\'egager 50\,\% et 100\,\% des axes primaires/secondaires.

\subsection{Sc\'enario 2 --- Mobilit\'e \'economique}

\noindent\textbf{Argumentaire.}
Les pertes \'economiques dues aux chutes de neige repr\'esentent plusieurs
centaines de M\$ par jour au Qu\'ebec~\cite{statcan}. Le r\'eseau commercial
montr\'ealais, concentr\'e sur les art\`eres principales, g\'en\`ere l'essentiel
de l'activit\'e \'economique. D\'egager les axes commerciaux et les couloirs de
transport en commun en priorit\'e limite l'impact sur les commerces.

\noindent\textbf{D\'efinition des priorit\'es :}
\begin{itemize}
  \item Priorit\'e 3 (haute) : \texttt{primary}, \texttt{secondary}
    (axes commerciaux et de transit)
  \item Priorit\'e 2 (moyenne) : \texttt{tertiary} (voies desservant les zones
    d'emploi)
  \item Priorit\'e 1 (basse) : voies locales r\'esidentielles
\end{itemize}

\noindent\textbf{B\'en\'efices attendus :}
r\'eduction des pertes \'economiques ; maintien du transport en commun.
\textbf{Cibles :} travailleurs, commer\c{c}ants, usagers des bus.
\textbf{Risques :} in\'egalit\'e g\'eographique (quartiers r\'esidentiels
\'eloign\'es d\'elaiss\'es).

\noindent\textbf{Indicateurs sp\'ecifiques :}
temps pour d\'egager 80\,\% des axes commerciaux ;
nombre d'arr\^ets de bus accessibles apr\`es 3\,h.

\subsection{Sc\'enario 3 --- \'Equit\'e r\'esidentielle}

\noindent\textbf{Argumentaire.}
Le rapport Lowrie (2025)~\cite{lowrie} souligne que les quartiers r\'esidentiels
sont syst\'ematiquement d\'eneig\'es apr\`es les grands axes, laissant les
pi\'etons (enfants, personnes \^ag\'ees, PMR) dans des conditions dangereuses.
L'INSPQ~\cite{inspeq} recommande de prioriser les rues r\'esidentielles pr\`es
des \'ecoles et CHSLD pour r\'eduire les accidents hivernaux de $-15\,\%$.

\noindent\textbf{D\'efinition des priorit\'es (ordre invers\'e) :}
\begin{itemize}
  \item Priorit\'e 3 (haute) : \texttt{residential}, \texttt{living\_street}
  \item Priorit\'e 2 (moyenne) : \texttt{tertiary}, \texttt{unclassified},
    \texttt{service}
  \item Priorit\'e 1 (basse) : \texttt{primary}, \texttt{secondary}
    (couvertes prioritairement par le sel)
\end{itemize}

\noindent\textbf{B\'en\'efices attendus :}
r\'eduction des accidents pi\'etons hivernaux ;
accessibilit\'e des \'ecoles et garderies.
\textbf{Cibles :} r\'esidents pi\'etons, personnes \^ag\'ees, enfants.
\textbf{Risques :} d\'elai accru pour les art\`eres commerciales et urgences.

\noindent\textbf{Indicateurs sp\'ecifiques :}
temps pour d\'egager 50\,\% et 100\,\% des rues r\'esidentielles ;
nombre de quartiers r\'esidentiels couverts apr\`es 6\,h.

%=============================================================
\section{Analyse des r\'esultats}
%=============================================================

\subsection{R\'esultats par arrondissement}

\begin{center}
\textbf{Tableau 2} --- R\'esultats du PPC-D\\[0.5em]
\begin{tabular}{|l|r|r|r|r|}
\hline
\textbf{Arrondissement} & \textbf{Circuit (km)} & \textbf{D. vide (km)} & \textbf{Ratio (\%)} & \textbf{Couv. (\%)} \\
\hline
Outremont              &   80,5 & 11,3 & 14,1 & 97,3 \\
Verdun                 &  110,5 &  8,6 &  7,8 & 99,1 \\
Anjou                  &  263,3 & 44,4 & 16,9 & 98,4 \\
Rivi\`ere-des-Prairies--PAT &  775,8 & 79,0 & 10,2 & 99,5 \\
\hline
\end{tabular}
\end{center}

\medskip
\begin{center}
\textbf{Tableau 3} --- Co\^uts op\'erationnels (3 v\'ehicules par arrondissement)\\[0.5em]
\begin{tabular}{|l|r|r|r|}
\hline
\textbf{Arrondissement} & \textbf{Distance/v\'eh. (km)} & \textbf{Dur\'ee/v\'eh. (h)} & \textbf{Co\^ut total (\$)} \\
\hline
Outremont              &  26,8 &  2,7 &  1\,597 \\
Verdun                 &  36,8 &  3,7 &  1\,634 \\
Anjou                  &  87,8 &  8,8 &  1\,819 \\
Rivi\`ere-des-Prairies--PAT & 258,6 & 25,9 &  2\,449 \\
\hline
\textbf{Total 4 arrondissements} & & & \textbf{7\,499} \\
\hline
\end{tabular}
\end{center}

\noindent\textbf{Note :} \`A Anjou ($0{,}8$\,h sup.) et RDP-PAT ($17{,}9$\,h sup.),
les v\'ehicules d\'epassent les 8\,h normales. Pour RDP-PAT, il est fortement
recommand\'e d'augmenter la flotte \`a 10 v\'ehicules, ce qui r\'eduit le temps
\`a $\approx 7{,}8$\,h/v\'ehicule (sans heures suppl\'ementaires) pour un
co\^ut total de $\approx 9\,500$\,\$.

\medskip
\noindent\textbf{Mod\`ele de co\^ut global.}
Pour $n$ v\'ehicules se partageant un circuit de $D$\,km :
\begin{equation}
  C(n) = n \cdot \left[500 + 1{,}1 \cdot \frac{D}{n}
  + f\!\left(\frac{D}{10n}\right)\right]
  \quad\text{o\`u}\quad
  f(t) = 1{,}1\,\min(t,8) + 1{,}3\,\max(t-8,0)
\end{equation}
Le co\^ut marginal d'un v\'ehicule suppl\'ementaire vaut $500$\,\$ (fixe) moins
l'\'economie r\'ealis\'ee sur les heures suppl\'ementaires --- profitable d\`es
que la dur\'ee/v\'ehicule d\'epasse 8\,h.

\subsection{Comparaison des sc\'enarios}

\begin{center}
\textbf{Tableau 4} --- Temps (h) pour d\'egager 50\,\% et 80\,\% des arcs haute priorit\'e\\[0.5em]
\begin{tabular}{|l|l|r|r|r|r|}
\hline
\textbf{Sc\'enario} & \textbf{M\'etrique} & \textbf{Outremont} & \textbf{Verdun} & \textbf{Anjou} & \textbf{RDP--PAT} \\
\hline
S1 -- Urgences     & 50\,\% HP (h) &  4,7 &  5,2 & 17,4 & 48,2 \\
                   & 80\,\% HP (h) &  7,2 &  8,7 & 21,6 & 67,9 \\
\hline
S2 -- \'Economique & 50\,\% HP (h) &  4,7 &  5,2 & 17,4 & 48,2 \\
                   & 80\,\% HP (h) &  7,2 &  8,7 & 21,6 & 67,9 \\
\hline
S3 -- R\'esidentiel & 50\,\% HP (h) &  4,0 &  3,9 &  9,2 & 37,8 \\
                   & 80\,\% HP (h) &  6,3 &  8,8 & 14,8 & 59,6 \\
\hline
\end{tabular}
\end{center}

\subsection{Projection des effets sur les habitants}

\noindent\textbf{Sc\'enario 1 (Urgences).}
Les art\`eres primaires et secondaires ($\approx 15\,\%$ des arcs) assurent
l'acc\`es \`a tous les services d'urgence. Elles sont d\'egag\'ees \`a 100\,\%
\`a Outremont en $\approx 8$\,h avec 3 v\'ehicules. \`A Anjou et RDP-PAT,
les temps d\'epassent 17\,h pour 50\,\% des axes prioritaires, ce qui justifie
d'augmenter la flotte dans les grands arrondissements.

\noindent\textbf{Sc\'enario 2 (\'Economique).}
Les indicateurs sont identiques \`a S1 car les deux sc\'enarios utilisent la
m\^eme classification \texttt{highway}. Une am\'elioration serait d'utiliser
des buffers g\'eospatiaux autour des zones commerciales (donn\'ees POI OSM) pour
diff\'erencier S2 de S1. La principale valeur de S2 est donc conceptuelle :
justifier la priorit\'e par l'impact \'economique plut\^ot que par la s\'ecurit\'e.

\noindent\textbf{Sc\'enario 3 (\'Equit\'e r\'esidentielle).}
Les rues r\'esidentielles constituent 65--75\,\% du r\'eseau. S3 d\'egage
50\,\% des rues r\'esidentielles \textbf{15--25\,\% plus t\^ot} que S1/S2
\`a Outremont et Verdun (4,0\,h vs 4,7\,h et 3,9\,h vs 5,2\,h). Ce gain
est significatif pour les pi\'etons : chaque heure \'economis\'ee r\'eduit
l'exposition au risque de chute sur la glace.
Revers de la m\'edaille : les art\`eres sont d\'egag\'ees plus tard, ce qui
peut g\^ener les navetteurs si les chasse-neige ne travaillent pas en parall\`ele
avec les camions de sel.

\noindent\textbf{Critique des sc\'enarios.}
\begin{itemize}
  \item S1 et S2 sont identiques dans l'impl\'ementation actuelle ;
    des donn\'ees POI permettraient de les diff\'erencier.
  \item S3 am\'eliore l'exp\'erience des r\'esidents pi\'etons au prix d'un l\'eger
    retard sur les art\`eres, acceptable si celles-ci ont \'et\'e trait\'ees au sel.
  \item Pour les grands arrondissements (Anjou, RDP-PAT), le mod\`ele \`a un seul
    v\'ehicule est irr\'ealiste ; la priorit\'e devrait s'exprimer dans un
    d\'ecoupage en zones (VRP) o\`u chaque v\'ehicule a sa propre priorit\'e
    de quartier.
\end{itemize}

\begin{thebibliography}{9}
\bibitem{osmnx}
  B. Boeing,
  \textit{OSMnx: New Methods for Acquiring, Constructing, Analyzing, and
  Visualizing Complex Street Networks},
  Computers, Environment and Urban Systems, vol.~65, 2017.
\bibitem{statcan}
  Statistique Canada,
  \textit{Impact \'economique des intemp\'eries hivernales},
  catalogue no~71-607, 2022.
\bibitem{lowrie}
  M. Lowrie,
  \textit{Comment \c{c}a marche, le d\'eneigement et le d\'egla\c{c}age
  \`a Montr\'eal ?},
  Le Devoir, d\'ecembre 2025.
\bibitem{sante}
  Sant\'e Montr\'eal,
  \textit{Urgences et conditions hivernales : retours d'exp\'erience},
  rapport interne, 2019.
\bibitem{inspeq}
  INSPQ,
  \textit{D\'eneigement et s\'ecurit\'e des pi\'etons : recommandations},
  Institut national de sant\'e publique du Qu\'ebec, 2021.
\end{thebibliography}

\end{document}
